
In the background Chapters, the authors should provide a concise, but self-contained theoretical exposition of the main technical tools to be used in the thesis. In principle, the technical content provided in the background section should enable the reader to understand the exposition of the main results.

QUESTION 1: One recommendation is to split the background part of the thesis into chapters or sections, each corresponding to a key topic covered in the thesis. What are the main topics covered in the background chapters/sections? 

QUESTION 2: Are the topics covered in the background chapters/sections of relevance for addressing the Researach Question?  
ANSWER: (yes/somewhat/no)

QUESTION 3: The exposition used in the background chapters/sections should be just be based in the literature. Especially, claims of a technical nature. Are the authors making technical claims that are not supported by the literature? 
ANSWER: (yes/no) If yes, the authors should provide references. 

QUESTION 4: The exposition on each background chapter/section should should have a good balance between text and technical content (formulas, algorithms, pseudocode, diagrams, etc). Too much text may indicate that there is a lack of depth in the technical part. Way too much technical content without connecting text may make the chapter difficult to read. Is the distribution between textual information and technical content well balanced? 
ANSWER: (well balanced/ way too much text/ technical content is way too dense) 

QUESTION 5: Some students may feel compelled to fill space by writing a lot of repeated information. Circular argumentation, reversion of the flow of ideas etc, are some common techniques. Are these unnecessary strategies for filling text being used? 
ANSWER: (yes/no) If yes, the student should shorten the excess of text by writing it as concisely as possible. 

QUESTION 6: Another strategy students use to fill in space is to use unnecessarily large figures. Does the chapter contains such figures? Ideally, a figures should be as compact as possible, preserving readability. 
ANSWER: (yes/no) If yes, students should scale the figures to a reasonable size. 


